\documentclass[a4paper,11pt]{article}

\usepackage[ngerman]{babel}
\usepackage[top=3cm,bottom=3cm,left=3cm,right=3cm]{geometry} %NICHT IN KEMIE2


%\usepackage{microtype} %schöneres typesetting  mit [stretch=10] für nur 1% anstelle von 2%
%\usepackage{lmodern} %Latin Modern FONT
%\usepackage{pifont} %schönere Font



\usepackage{amsmath}
\usepackage{nicefrac}
\usepackage{amssymb}
%\usepackage{mathtools}

\usepackage{titlesec} %Um Section, etc bearbeiten zu können
\usepackage{setspace}
\usepackage{enumitem} %enumerate
\usepackage{multicol}
\usepackage[many]{tcolorbox}
\usepackage{hyperref}
\usepackage{arydshln}

%\renewcommand{\ttdefault}{lmtt} %Schriftart wird geändert zu Latin Modern Typewriter


\hypersetup{hidelinks,
			colorlinks=true,
			linkcolor=black,
			citecolor=miudiutex,
			urlcolor=miugray2
			}

\setlength{\parindent}{0pt}
\setcounter{section}{0}


\titlespacing*{\section}{0pt}{20mm}{6mm}
\titlespacing*{\subsection}{0pt}{6mm}{3mm}
\titlespacing*{\subsubsection}{0pt}{3mm}{0mm}



%\definecolor{miudiutex}{HTML}{2f3d39}
\definecolor{miudiutex}{HTML}{1d2926}
\definecolor{miugray}{HTML}{b4cccf}
\definecolor{miugray2}{HTML}{4d7365}
\definecolor{red}{HTML}{cf1f5c}
\definecolor{green}{HTML}{00A36C}
\definecolor{delphinidin}{HTML}{315794}
\definecolor{pelargonidin}{HTML}{b33807}
\definecolor{cyanidin}{HTML}{ba0746}
\definecolor{petunidin}{HTML}{b56fed}



\raggedcolumns %wichtig für Multicolumns


%================================
% SYMBOLS
%================================

\newcommand{\RA}{$\rightarrow$~}
\newcommand{\RL}{$\rightleftharpoons$~}
\newcommand{\HARP}{\ding{226}~}
%$\xrightarrow{2}$ %Pfeil mit 2 drüber


%================================
% SHORT COMMANDS (BOXES ; ETC)
%================================


\newcommand{\notiz}[1]{\begin{notizz}{}{}\begin{center}
			{\color{miugray2!50!miudiutex}#1} \end{center}\end{notizz}}
\newcommand{\zit}[2]{\begin{zitat*}{#1}{}\small #2 \end{zitat*}}

\newcommand{\frage}[2]{\begin{qst*}{#1}{}\small #2 \end{qst*}}

\newcommand{\unten}{\end{center} \tcblower \begin{center}}


\newcommand*\circled[1]{\tikz[baseline=(char.base)]{
            \node[shape=circle,draw,inner sep=0.5pt,miudiutex] (char) {\tiny #1};}}
            
\newcommand{\miusquare}{{\color{miugray2!50}\scriptsize $\blacksquare$}~}
\newcommand{\miusquarp}{{\color{petunidin!50}\scriptsize $\blacktriangleright$}~}
\newcommand{\niu}{\item[\miusquare]}
\newcommand{\nip}{\item[\miusquarp]}

%%%%% ABSAETZE

\newcommand{\pps}{\par \vspace*{1mm}}
\newcommand{\pp}{\par \vspace*{3mm}}
\newcommand{\ppbig}{\par \vspace*{8mm}}
\newcommand{\pphuge}{\par \vspace*{10mm}}

%================================
% FRAGE BOX
%================================

\tcbuselibrary{theorems,skins,hooks}
\newtcbtheorem{qst}{Frage:}
{%
	enhanced,
	breakable,
	colback = gray!15!white,
	frame hidden,
	boxrule = 0sp,
	borderline west = {2.5pt}{0pt}{red!70},
	sharp corners,
	detach title,
	before upper = \tcbtitle\par\smallskip,
	coltitle = black,
	fonttitle = \bfseries\sffamily,
	description font = \mdseries,
	separator sign none,
	segmentation style={solid, gray!50!black},
}
{th}




%================================
% ZITAT BOX
%================================

\tcbuselibrary{theorems,skins,hooks}
\newtcbtheorem{zitat}{Zitat}
{%
	enhanced,
	breakable,
	colback = gray!15!white,
	frame hidden,
	boxrule = 0sp,
	borderline west = {2.5pt}{0pt}{black!70},
	sharp corners,
	detach title,
%	before upper = \tcbtitle\par\smallskip,
	coltitle = gray!50!black,
	fonttitle = \bfseries\sffamily,
	description font = \mdseries,
%	separator sign none,
	segmentation style={solid, gray!50!black},
}
{th}



%================================
% SIMPLE SCHWARZER RAHMEN BOX
%================================


\newcommand{\boxfr}[1]{
	\begin{tcolorbox}[
	drop shadow southeast,
	enhanced,
	colframe=black!50,
	colback=white,
	boxrule = 1pt, 
%	toprule=0pt, bottomrule=0pt,rightrule=0pt,leftrule=0pt,
	boxsep=5pt,
	arc=1pt,
	]
	\begin{center}
	#1
	\end{center}
	\end{tcolorbox}
}

\definecolor{miudiutex}{HTML}{1d2926}
\definecolor{miugray}{HTML}{b4cccf}
\definecolor{miugray2}{HTML}{4d7365}
\definecolor{white2}{HTML}{f5f7f8}
\definecolor{red}{HTML}{cf1f5c}
\definecolor{green}{HTML}{00A36C}
\definecolor{delphinidin}{HTML}{315794}
\definecolor{uniblau}{HTML}{004e9f}
\definecolor{unigelb}{HTML}{fcba00}
\definecolor{unigrau}{HTML}{909085}
\definecolor{pelargonidin}{HTML}{b33807}
\definecolor{cyanidin}{HTML}{ba0746}
\definecolor{paeonidin}{HTML}{d15ca6}
\definecolor{petunidin}{HTML}{b56fed}
\definecolor{malvidin}{HTML}{8a2d8c}

\newcommand{\metermilli}{\mskip3mu \text{mm}}
\newcommand{\cm}{\mskip3mu \text{cm}}
\newcommand{\m}{ \text{m}}
\newcommand{\lite}{ \text{l}}
\newcommand{\kg}{ \text{kg}}
\newcommand{\g}{ \text{g}}
\newcommand{\km}{ \text{km}}
\newcommand{\h}{ \text{h}}
\newcommand{\s}{ \text{s}}
\newcommand{\tonne}{ \text{t}}
\usepackage[utf8]{inputenc}
\usepackage[T1]{fontenc}
\usepackage{lmodern}

\usepackage[
    activate={true,nocompatibility},
    final,
    tracking=true,
    kerning=true,
    spacing=true,
    factor=1100,
    stretch=30,
    shrink=20
]{microtype}
\usepackage{pdfpages}
\usepackage{awesomebox}
\usepackage{marginnote}
\tcbuselibrary{documentation}
\usepackage{sidenotes}
\usepackage{import}
\usepackage{xifthen}
\usepackage{transparent}
\usepackage[table]{xcolor}


\newcommand{\abbicon}{%
  \protect\tikz[baseline=0.2mm,scale=0.8]{%
    % Das blaue Quadrat mit abgerundeten Ecken
    \fill[uniblau, rounded corners=1.2pt] (0,0) rectangle (0.8em, 0.8em);
    % Der weiße Kreis in der Mitte
    \fill[white] (0.4em, 0.4em) circle (0.12em);
  }%
} 
\usepackage[labelfont={bf},font={color=miudiutex,small},figurename={\abbicon $~$Abbildung}]{caption}


\newcommand{\bckgrnd}{
\AddToHookNext{shipout/background}{
\begin{tikzpicture}[remember picture, overlay]
 \draw[draw=none,fill=uniblau!70, shift={(current page.north west)}] (-3,0) rectangle (23,-3.75);
 \draw[draw=none,fill=miugray, shift={(current page.north east)}] (-7,0) rectangle (3,-3.75);
\end{tikzpicture}
}
}



\newcommand{\incfig}[2]{%
\def\svgwidth{#2\columnwidth}
\import{C://LaTeX-Files/pngs/}{#1.pdf_tex}
}

\renewcommand{\textbf}[1]{\textsf{{\bfseries {#1}}}}



\titleformat{\section}[hang]
		{\LARGE \color{uniblau!75!black} \sffamily \bfseries}
		{\thesection}
		{6mm}
		{}
		[]
\titleformat{\subsection}[hang]
		{\large \sffamily \bfseries \color{black}}
		{\thesubsection}
		{4mm}
		{{\color{miugray}$\blacksquare ~$}}

\titleformat{\subsubsection}[block]
		{\normalsize \bfseries \sffamily \color{miudiutex}}%
		{\normalsize \color{black} \thesubsubsection}
		{2mm}
		{{\color{miugray}$\blacksquare ~$}}
		[ \color{black}]
		
\titleformat{\paragraph}
		{\normalsize \bfseries \sffamily \color{black}}
		{\footnotesize \theparagraph}
		{1mm}
		{}

\pagestyle{empty}
\titlespacing*{\section}{0pt}{20mm}{12mm}
\titlespacing{\section}{0pt}{20mm}{12mm}

\titlespacing*{\subsection}{0pt}{14mm}{4mm}
\titlespacing{\subsection}{0pt}{14mm}{4mm}

\titlespacing*{\subsubsection}{0pt}{6mm}{3mm}
\titlespacing{\subsubsection}{0pt}{6mm}{3mm}

\titlespacing*{\paragraph}{0pt}{6mm}{2mm}
\titlespacing{\paragraph}{0pt}{6mm}{2mm}


\let\oldmarginfigure\marginfigure
\let\endoldmarginfigure\endmarginfigure

\let\oldmarginfigure\marginfigure
\let\endoldmarginfigure\endmarginfigure

\renewenvironment{marginfigure}[1][0pt] % [1] steht für 1 Argument, [0pt] ist der Standardwert
  {\oldmarginfigure[#1]\captionsetup{labelfont={sf,bf,color=uniblau!75!black},font={color=miudiutex,small}}}
  {\endoldmarginfigure}
  
  
  
  
\renewcommand{\labelitemi}{\color{uniblau!70}$\bullet$}
\renewcommand{\labelitemii}{\color{uniblau!70}$\cdot$}
\newcommand{\plmtsquare}{{\color{uniblau!70}$\blacksquare~$}}
\newcommand{\splmt}[1]{{\color{uniblau!75!black} \sffamily #1}}
\newcommand{\fplmt}[1]{{\color{uniblau!75!black} \sffamily \bfseries #1}}

\newcommand{\antwort}[1]{
\begin{tcolorbox}[
	enhanced,
	enforce breakable,
	standard jigsaw,
	boxrule=0.7pt,
	colframe=black,
	sharp corners,
	boxsep=5pt,
	title=\bfseries \sffamily Antwort,
	opacityback=0,
	coltitle=miudiutex,
	attach boxed title to top text left={yshift=-\tcboxedtitleheight/2},
	boxed title style={colback=white,colframe=white},
	bottomrule at break=0pt,
	toprule at break=0pt,
	pad at break=16mm,
]
	\begin{center}
	#1
	\end{center}
	\end{tcolorbox}
}




\rowcolors{2}{miugray!50}{miugray!50}
\arrayrulecolor{white} % Weiße Gitterlinien
\setlength{\arrayrulewidth}{1.5pt} % Etwas dickere Linien
\renewcommand{\arraystretch}{1.4} % Mehr Zeilenhöhe für bessere Lesbarkeit}


%\setlist[itemize]{itemsep=0mm}
\setlist[enumerate]{itemsep=0mm}


\newif\ifshowme
%\showmetrue




\begin{document}
\newgeometry{top=1.8cm,bottom=4cm,left=1.5cm,right=7.5cm,marginparsep=-2.00cm,marginparwidth=4.75cm}


\bckgrnd


\begin{center}
\LARGE
\color{white}
\textbf{Übung 03} \par \vspace*{2mm}
\large \color{miugray} \textbf{Prozessbezogene Lebensmitteltechnologie}
\end{center}
\par 
\vspace*{6mm}


\color{miudiutex}
\section*{Wärmeübertragung I}
\subsection*{Aufgabe 01}
Berechnen Sie die Strahlungsleistung, die von 100 m$^2$ einer polierten Eisenoberfläche
(Emissionsgrad = 0,06) abgegeben wird. Die Temperatur der Oberfläche beträgt 37 °C.

%\begin{marginfigure}
%\centering
%\includegraphics[width=4.5cm]{F://LaTeX-Files/pngs/strawberry}
%\caption{Erdbeerbecher mit Vanilleschmand und Gewürzcrumble}
%\end{marginfigure}

\ifshowme
\antwort{
\begin{align*}
P &= \varepsilon \cdot \sigma \cdot A \cdot T^4 \\ 
P &= 0,06 \cdot  100~\text{m}^2 \cdot 5,67 \cdot 10^{-8} ~\text{W} \cdot \text{m}^{-2}\text{K}^{-4}\cdot 310,15~\text{K}^4 \\
P &= 3.147,9~\text{W}
\end{align*}
}
\else
\fi

\subsection*{Aufgabe 02}
Eine Seite einer 1 cm dicken Edelstahlplatte wird auf 110 °C gehalten, die andere Seite
auf 90 °C. Berechnen Sie unter der Annahme stationärer Bedingungen die
Wärmeübertragungsrate pro Flächeneinheit durch die Platte. Die Wärmeleitfähigkeit von
nichtrostendem Stahl beträgt \mbox{17 W m$^{-1}$ °C$^{-1}$}.

\ifshowme
\antwort{
\begin{align*}
\dot{Q} &= - A \cdot \frac{\lambda}{d} \cdot (T_2 - T_1) \\
- \frac{\dot{Q}}{A} &= \cdot \frac{\lambda}{d} \cdot (T_2 - T_1) \\
- \frac{\dot{Q}}{A} &= \cdot \frac{17 \text{W} / \text{m°C}}{0,01m} \cdot 20\text{°C}\\
	&= 34.000 ~\text{W / m$^2$ s}
\end{align*}

}
\pagebreak
\else
\fi

\subsection*{Aufgabe 03}
Ein Glühofen hat eine 4 m$^2$ große Wand aus \splmt{Schamotte}. Die Dicke der Wand sei 0,5 m.
Zwischen der Innen- und der Außenseite der Wand besteht eine Temperaturdifferenz von
1.000 K. Die Temperatur außerhalb des Ofens beträgt 20 °C. Die Wärmeleitfähigkeit der
Schamotte ist $\lambda_1$ = 0,93 W/mK. 


\marginnote[]{\small Als \fplmt{Schamotte} werden im allgemeinen Sprachgebrauch häufig alle feuerfesten Steine und Ausmauerungen bezeichnet.}[-2.5cm]

\begin{enumerate}[label=\alph*)]
\item Wie groß ist der tägliche Wärmeverlust durch die Wand?


\ifshowme
\antwort{
\begin{align*}
\dot{Q} &= - A \cdot \frac{\lambda}{d} \cdot (T_2 - T_1) \\
&= - 4m^2 \cdot \frac{0,93}{0,5m} \cdot 1000 K \\
&= - 7.440~W \\
&= - 7.440~W \cdot 24h \\
&= - 178.560~Wh
\end{align*}
}
\else
\fi

\item Um wie viel Prozent verringert sich dieser Verlust, wenn die Außenwand (von außen)
durch eine 3 cm dicke Isolierschicht mit einer Wärmeleitfähigkeit von $\lambda_2$ = 0,093 
W/mK verstärkt wird?


\ifshowme
\antwort{
\begin{align*}
\dot{Q} &= - \frac{A \cdot \Delta T_{ges}}{(\frac{d_2}{\lambda_2}+\frac{d_1}{\lambda_1})} \\
\end{align*}
Einsetzen:
\begin{align*}
\dot{Q} &= - \frac{4m^2 \cdot 1000 K}{(\frac{0,03m}{0,093 W/mk}+\frac{0,5m}{0,93 W/mk})} \\
&= 4.650 W
\end{align*}
Prozentualer Unterschied:
\begin{align*}
\frac{4.650 W}{7.440 W} &\approx  0,625 = 62,5\%
\end{align*}
Also verringert sich der Verlust um 100-62,5\% = 37,5\%

}
\pagebreak
\else
\fi

\item Welche Temperatur herrscht an der Grenze zwischen der Isolierung und der
Schamotte-Schicht?

\ifshowme
\begin{marginfigure}[0cm]
\else
\begin{marginfigure}[-4cm]
\fi

\centering
\begin{tikzpicture}
\draw[fill=miugray!50] 	(0,0) rectangle (0.8,3);
\draw[fill=white]		(0.8,0) rectangle (2,3);
\node at (0.4,1.5) {$\dot{Q_1}$};
\node at (1.4,1.5) {$\dot{Q_2}$};
\draw[->,shorten >=2pt] 	(0,-0.3) -- (0,0)	node[below,yshift=-2mm]{$T_1$}; 
\draw[->,shorten >=2pt] 	(0.8,-0.3) -- (0.8,0)node[below,yshift=-2mm]{$T_2$};
\draw[->,shorten >=2pt] 	(2,-0.3) -- (2,0)	node[below,yshift=-2mm]{$T_3$};
\end{tikzpicture}
\caption{Der Wärmestrom ist an jeder Stelle der Wand konstant ($\dot{Q_1} = \dot{Q_2}$). Die Differenz der Temperatur $\Delta T$ lässt sich in ihre Bestandteile $T_1$, $T_2$ und $T_3$ zerlegen.}
\end{marginfigure}

\ifshowme
\antwort{
\begin{align*}
\dot{Q} &= - \lambda \cdot A \cdot \frac{T_3 - T_2}{d} \\
T_2 &= \frac{\dot{Q} \cdot d}{\lambda \cdot A } + T_3 \\
\\
T_2 &= \frac{4.650 ~\text{W} \cdot 0,03~\text{m}}{0,093 ~\text{W/mk} \cdot 4~\text{m}^2} + 20~^{\circ}\text{C} \\
 &= 395~^{\circ}\text{C}
\end{align*}
}

\else
\fi

\end{enumerate}


\subsection*{Aufgabe 4}
Welche Wärmemenge (in kWh) benötigen Sie, um 6,5 L einer Kartoffelsuppe von 10 °C auf
60 °C zu erhitzen? Es wird angenommen, dass die Dichte der Suppe 1 kg/dm$^3$ beträgt und
es keine Wärmeverluste gibt. \pp 
Die Suppe setzt sich pro 100 g zusammen aus: \par 
\begin{itemize}
\item Fett: 14,3 g
\item Proteine: 2,7 g
\item Kohlenhydrate: 11,5 g
\item Weitere Bestandteile: 0,5 g
\end{itemize}

%\boxfr{
%\begin{align*}
%c = &\quad ~   4,180 \times \text{Wasser} \\& + 1,711 \times \text{Proteine} \\& + 1,928 \times \text{Fett} \\&+ 1,547 \times \text{Kohlenhydrate} \\& + 0,908 \times \text{weitere Bestandteile}
%\end{align*}
%\pp
%}
\marginnote[]{\small Für die Berechnung der spezifischen \fplmt{Wärmekapazität $c$} kann folgende Formel verwendet
werden:

\begin{align*}
c = &\quad ~   4,180 \times \text{Wasser} \\& + 1,711 \times \text{Proteine} \\& + 1,928 \times \text{Fett} \\&+ 1,547 \times \text{Kohlenhydrate} \\& + 0,908 \times \text{weitere Bestandteile}
\end{align*}
}[-6cm]

\pp


\ifshowme
\antwort{
Zunächst Wärmekapazität $c$ ausrechnen:
\begin{align*}
c &= ~~~ 4,180 \cdot 0,71 + 1,711 \cdot 0,027 + 1,928 \cdot 0,143 \\
&~~~ + 1,547 \cdot 0,115 + 0,908 \cdot 0,005 \\
 &= ~~~ 3,472146~ \text{kJ} \cdot \text{kg}^{-1} \cdot \text{K}^{-1}
\end{align*}
Energie berechnen:
\begin{align*}
Q &= m \cdot c \cdot \Delta T \\ \\
Q &= 6,5~ \text{kg} \cdot 3,472~ \text{kJ} \cdot \text{kg}^{-1}\cdot \text{K}^{-1} \cdot (333,15-283,15 ~ \text{K}) \\
&= 1.128,45~ \text{kJ} \\
&= 1.128,45~ \text{kWs}  \\
&= 0,313~ \text{kWh}
\end{align*}

}
\pagebreak
\else
\fi



\subsection*{Aufgabe 5}
In Ihrem 20°C warmen Zuhause schalten Sie Ihre Herdplatte auf maximaler Stufe an. Diese erhitzt sich daraufhin auf 300°C. Im nächsten Schritt stellen Sie eine Pfanne aus Aluminium ($\lambda_{Aluminium} = 236~\text{W/mK}$) mit einem Durchmesser von 28 cm und einer Dicke von 2 mm auf die Herdplatte.
\begin{itemize}
\item Welcher Wärmefluss entsteht zwischen der Herdplatte und der Innenseite der Pfanne?
\ifshowme
\antwort{
\begin{align*}
\dot{Q} &= \lambda_{Al} \cdot \frac{A}{d} \cdot (\Delta T) \\
&= 236~\text{W/mK} \cdot \frac{\pi \cdot 0,15~\text{m}^2}{0,002~\text{m}} \cdot 280~\text{K}\\
&= 4.956.000~\text{J}\\
&= 4.956~\text{kJ}
\end{align*}
}
\else
\fi

\item Wieso können Sie von dem fließenden Wärmestrom keine Rückschlüsse auf ihre tatsächliche Stromrechnung schließen?
\ifshowme
\antwort{
Das Delta T wird mit der Zeit immer geringer. Die Herdplatte muss anfänglich eine viel größere Leistung aufbringen um die kalte Herdplatte (bzw auch die Pfanne) zu erhitzen.
}
\else
\fi

\end{itemize}




\end{document}
